% Ppopy Phase 2 Submission A
% Revised from the submitted PDF and aligned with the Phase 2 notice.
% This file contains the complete document; IEEEtran.cls is included alongside it.
\documentclass[conference,letterpaper]{IEEEtran}
\usepackage[T1]{fontenc}
\usepackage{courier}
\usepackage{amsmath,amssymb}
\usepackage{booktabs,array,tabularx,cite,url,listings,xcolor,tikz}
\usepackage[hidelinks]{hyperref}
\urlstyle{same}
\hypersetup{pdftitle={Ppopy: Proactive Skincare Decision Support with a Smart Mirror},pdfauthor={Doyeong Koo; Takgeon Kang; Geonil Won}}
\lstset{basicstyle=\ttfamily\footnotesize,breaklines=true,columns=fullflexible,keepspaces=true,showstringspaces=false,frame=single,rulecolor=\color{gray},aboveskip=5pt,belowskip=5pt,xleftmargin=2pt,xrightmargin=2pt}
\title{Ppopy: Proactive Skincare Decision Support with a Smart Mirror}
\author{\IEEEauthorblockN{Doyeong Koo, Takgeon Kang, Geonil Won}\IEEEauthorblockA{Hanyang University\\Team CosMate -- CSE406 Software Engineering}}
\begin{document}
\maketitle
\pagestyle{plain}
\thispagestyle{plain}
\begin{abstract}
Team CosMate proposes ppopy to help users decide what to maintain, adjust or check before their skincare routine. The system compares confirmed product information, planned routines and reported use with a knowledge base (KB), and requests only information that can change a decision. For proactive guidance, it estimates the next routine start from user-confirmed use times, selects any needed guidance and displays it on a previously authorized smart mirror before that start. The mirror provides an overview at the washbasin or vanity; users answer questions and approve routine changes in the app. The prototype combines context-sensitive decisions with advance delivery on one home display.
\end{abstract}
\begin{IEEEkeywords}skincare decision support, knowledge base, proactive guidance, smart mirror\end{IEEEkeywords}
\section*{Role Assignments}
{\centering\footnotesize
\renewcommand{\arraystretch}{1.12}\setlength{\tabcolsep}{3pt}
\begin{tabularx}{\columnwidth}{@{}>{\hsize=.80\hsize\raggedright\arraybackslash}X>{\hsize=.70\hsize\raggedright\arraybackslash}X>{\hsize=1.50\hsize\raggedright\arraybackslash}X@{}}\toprule
\textbf{Roles} & \textbf{Name} & \textbf{Task description}\\\midrule
User & Takgeon Kang, Geonil Won & Review routine workflows, app questions and mirror usability.\\
Customer & Doyeong Koo & Define requirements, feature priorities, budget and acceptance criteria.\\
Software developer & All members & Koo: app/backend, KB and integration; Kang: agent questions and guidance; Won: mirror hardware and display. All: data collection and testing.\\
Development manager & All members (joint) & Share schedule and progress updates, agree on changes and review system integration.\\
\bottomrule
\end{tabularx}
}

\section{Introduction}

\subsection{Problem and Objective}

The project addresses a practical need: deciding how to carry out a skincare routine and what to check before starting it. A product list alone does not explain which part of a user's planned routine needs attention. ppopy combines confirmed products, actual-use reports, user goals and KB criteria to select relevant guidance. A mirror at the washbasin or vanity makes that guidance available while the user prepares for the routine.

\subsection{User Workflow and Example}

Users register products by name, photo or screenshot, confirm the product and variant, and record their plan and actual use. They set a goal such as ``show what I need to check before starting.'' The system reviews the plan and context, asks decision-changing questions in the app, and updates its findings after answers.

For example, three confirmed evening starts at 20:56, 21:00 and 21:04 produce an estimated start of 21:00. At 20:50, the system reviews the upcoming routine. It may select an adjustment, request a missing fact, or find no need for new guidance. With valid prior display approval, needed guidance appears on the mirror before 21:00. The estimated time determines when to review; the KB and user context determine what to do.

\subsection{Related Software}

SkinSort's routine creator supports routine organization \cite{ref1}; INKEEDecoder explains ingredient lists \cite{ref2}. These functions inform ppopy's product and routine presentation. The proposed workflow connects use-context comparison, decision-dependent questions and advance guidance on a home display.

\subsection{Prototype Scope and Current Status}

Development work includes product-input components, a local React/Python KB review, follow-up questions and account-separated JSON records \cite{ref3}. The course prototype will integrate these components with history-based timing, automatic delivery, real account authentication, cloud storage and one output-only mirror. Questions, settings and plan approval remain in the app. Time and weather form the mirror's default screen; broader appliance orchestration is future work. The timing estimate concerns the next routine start, not future skin condition.

Course code and the A/B LaTeX sources are available at \url{https://github.com/rnehdud/ppopy}.

\section{Requirements}

\subsection{Functional Requirements}

\subsubsection*{R1. Sign-in and consent}

Require sign-in and appropriate consent for personal routine records; restrict each record to its owner.

\subsubsection*{R2. Product registration and ingredients}

Confirm products and variants from names, photos or screenshots. Verify ingredient sources, product versions and standard-name matches before using them as facts.

\subsubsection*{R3. Plans, actual use and user intent}

Keep confirmed plans separate from actual-use reports. Allow edits to the time, products, use details and reported discomfort, and let users set guidance goals and preferences.

\subsubsection*{R4. Contextual KB review}

Review current and planned routines against applicable KB criteria and exceptions. Distinguish confirmed facts from unknown or conflicting information.

\subsubsection*{R5. Decision-dependent questions}

Ask only for unknown information that can change a decision. Update related findings after answers; preserve skipped or declined answers as unknown without repeating a question in unchanged context.

\subsubsection*{R6. Decision and guidance selection}

Select maintain, adjust, check or defer with linked reasons and evidence. Preserve the decision when no new user-facing intervention is needed. A suggestion changes the confirmed plan only after explicit app acceptance.

\subsubsection*{R7. Separate records and outcomes}

Store plans, suggestions, actual-use reports, time estimates, decisions and delivery/display outcomes separately, with their revisions and origins.

\subsubsection*{R8. Mirror pairing and display approval}

Pair one output-only mirror. Configure its room, timezone, default screen, allowed display periods and approval expiry in the app.

\subsubsection*{R9. Next-time estimation and proactive guidance}

Estimate the next routine start from valid actual-use times and review the latest context beforehand. Display needed, current guidance before that start within prior approval; prevent duplicates and invalidate outdated decisions or commands.

\subsubsection*{R10. Account and display controls}

Provide logout, consent/display-approval withdrawal and account deletion. End personal access and expire displayed content when authority ends; distinguish a deletion request from completed cleanup.

% Start the quality requirements in the final right column, as in the Revised PDF.
\newpage
\subsection{Quality Requirements}

\subsubsection*{Q1. Personal information protection}

Protect transmission and storage. Authorize personal mirror content by account, paired device and valid display approval, without inferring who is viewing the mirror.

\subsubsection*{Q2. Reliable execution and accessibility}

Discard stale guidance, prevent duplicate interventions and record only acknowledged delivery stages. Show errors clearly and support readable Korean controls, enlarged text and VoiceOver.

\subsubsection*{Q3. Traceability and verification}

Link estimates and decisions to inputs, parameters, KB sources and execution outcomes. Map R1--R10/Q1--Q3 to reproducible tests and distinguish planned tests from executed evidence.
\begin{thebibliography}{00}
\bibitem{ref1} SkinSort, ``Routine creator.'' [Online]. Available: \url{https://skinsort.com/support/the-routine-creator}. Accessed: Nov. 6, 2026.
\bibitem{ref2} INKEEDecoder, ``About INKEEDecoder.'' [Online]. Available: \url{https://inkeedecoder.com/about}. Accessed: Nov. 6, 2026.
\bibitem{ref3} Team CosMate, ``ppopy / SkinAnalysis, project repository,'' 2026. [Online]. Available: \url{https://github.com/rnehdud/SkinAnalysis}. Accessed: Nov. 6, 2026.
\end{thebibliography}
\end{document}
